\documentclass[11pt]{article}
\usepackage[a4paper,margin=22mm]{geometry}
\usepackage{fontspec}
\setmainfont{DejaVu Serif}
\setsansfont{DejaVu Sans}
\usepackage{microtype}
\usepackage{xcolor}
\usepackage{hyperref}
\usepackage{enumitem}
\usepackage{parskip}
\definecolor{Accent}{HTML}{971D32}
\definecolor{Muted}{HTML}{666666}
\hypersetup{colorlinks=true,urlcolor=Accent,linkcolor=Accent}
\setlist{leftmargin=1.6em,itemsep=0.3em,topsep=0.35em}
\setcounter{tocdepth}{2}
\renewcommand{\contentsname}{Contents}
\newcommand{\sectionrule}{\par\vspace{-0.5em}\noindent\textcolor{Accent}{\rule{\linewidth}{0.6pt}}\par}
\begin{document}

\begin{center}
{\sffamily\bfseries\color{Accent} TOEFL WORD OF THE DAY\par}
\vspace{8mm}
{\fontsize{42}{48}\selectfont\bfseries ubiquitous\par}
\vspace{2mm}
{\Large\sffamily\color{Accent} /juːˈbɪkwətəs/\par}
\vspace{2mm}
{\small\sffamily Local audio phrase: ubiquitous\par}
{\small\sffamily Audio file: audios/ubiquitous.mp3\par}
\vspace{2mm}
{\large\sffamily\color{Muted} adjective\par}
\vspace{5mm}
{\sffamily present almost everywhere\par}
\vspace{6mm}
{\large Something ubiquitous can be seen, found, or experienced almost everywhere.\par}
\vspace{6mm}
\begin{minipage}{0.84\textwidth}
\itshape “Smartphones have become ubiquitous in modern society.”
\end{minipage}\par
\vspace{10mm}
\textbf{Date:} July 28th 2026\quad\textbullet\quad \textbf{Level:} B1--mid-B2 TOEFL
\vfill
Created and compiled by \textbf{Amir Kooshky}\par
\vspace{2mm}
\href{https://t.me/KooshkyTOEFL}{Telegram Channel}\quad\textbullet\quad
\href{https://www.instagram.com/kooshkytoefl}{Instagram}\quad\textbullet\quad
\href{https://t.me/KooshkyTOEFL\_pv}{Personal Telegram}
\end{center}
\newpage
\tableofcontents
\newpage

\section{Meanings and usage}
\sectionrule
\subsection{Meaning 1: Present almost everywhere}
\textbf{Part of speech:} adjective

\textbf{IPA:} /juːˈbɪkwətəs/

\textbf{Pronunciation phrase:} ubiquitous

\textbf{Local audio file:} audios/ubiquitous.mp3

\textbf{Register:} neutral to formal; common in academic, scientific, technological, and social discussions

\textbf{Definition:} present, found, or appearing almost everywhere, especially within a particular place, society, or situation

\textbf{In simpler words:} almost everywhere

\textbf{Typical pattern:} ubiquitous in, across, or throughout + place or area

\subsubsection{Natural examples}
\begin{enumerate}
  \item Smartphones have become ubiquitous in modern society.
  \item Plastic is now ubiquitous in the world's oceans.
  \item Coffee shops are ubiquitous throughout the city.
  \item The researchers found that the bacteria were ubiquitous in the local environment.
  \item Advertising has become ubiquitous on social media platforms.
  \item Security cameras are increasingly ubiquitous in public spaces.
  \item The plant is ubiquitous across the region because it can survive in many different climates.
\end{enumerate}

\subsubsection{Nuance and implication}
\textbf{Central nuance:} The word does not usually mean literally present in every possible place. It emphasizes that something is extremely widespread and difficult to avoid or overlook.

\textbf{Usual implication:} It often suggests that something has become a normal or noticeable part of an environment or society.

\textbf{Compared with widespread:} Widespread means existing in many places or affecting many people. Ubiquitous is stronger and suggests that something seems to be almost everywhere.

\subsubsection{Grammar notes}
\textbf{How to use it:} Use it before a noun or after a linking verb such as be, become, or seem.

\textbf{What can follow it:} When it comes after a linking verb, it is often followed by in, across, or throughout and then a place, field, or situation.

\textbf{Common preposition:} Use ubiquitous in for an environment or area, as in ubiquitous in modern life. Across and throughout are also common when emphasizing geographical spread.

\textbf{Position in a sentence:} You can say ubiquitous technology before a noun, or Technology has become ubiquitous after a linking verb.

\subsubsection{Useful patterns}
\textbf{Pattern:} ubiquitous + noun

\textbf{Use:} Use this before a noun to describe something that is present almost everywhere.

\textbf{Example:} The study examines the effects of ubiquitous digital technology.

\textbf{Pattern:} be or become ubiquitous

\textbf{Use:} Use this when something is or becomes extremely common in many places.

\textbf{Example:} Wireless internet has become ubiquitous on university campuses.

\textbf{Pattern:} ubiquitous in + place or situation

\textbf{Use:} Use this to identify the environment in which something is found almost everywhere.

\textbf{Example:} These symbols are ubiquitous in religious art.

\textbf{Pattern:} ubiquitous across or throughout + area

\textbf{Use:} Use this to emphasize that something is widely present over a large area.

\textbf{Example:} The species is ubiquitous throughout the coastal region.

\subsubsection{Common wrong usage}
\paragraph{Correction 1}
\textbf{Wrong:} Mobile phones are ubiquitous everywhere.

\textbf{Correct:} Mobile phones are ubiquitous.

\textbf{Why:} Ubiquitous already means almost everywhere, so everywhere is usually unnecessary.

\paragraph{Correction 2}
\textbf{Wrong:} This technology is an ubiquitous feature of modern life.

\textbf{Correct:} This technology is a ubiquitous feature of modern life.

\textbf{Why:} Use a because ubiquitous begins with the consonant sound /j/, like the first sound in university.

\paragraph{Correction 3}
\textbf{Wrong:} Online advertising has ubiquitous in modern life.

\textbf{Correct:} Online advertising has become ubiquitous in modern life.

\textbf{Why:} Ubiquitous is an adjective, so use a linking verb such as become before it.

\section{Word family}
\sectionrule
\subsection{ubiquity}
\textbf{Part of speech:} uncountable noun

\textbf{IPA:} /juːˈbɪkwəti/

\textbf{Definition:} the state of being present or found almost everywhere

\textbf{Relationship to the headword:} This noun names the state or quality described by ubiquitous.

\textbf{Grammar pattern:} the ubiquity of + noun

\textbf{Usage note:} It is formal and is usually used in the expression the ubiquity of something.

\subsubsection{Examples}
\begin{enumerate}
  \item The ubiquity of smartphones has changed how people communicate.
  \item Researchers are concerned about the ubiquity of plastic waste in marine environments.
\end{enumerate}

\subsection{ubiquitously}
\textbf{Part of speech:} adverb

\textbf{IPA:} /juːˈbɪkwətəsli/

\textbf{Definition:} in a way that is present or occurs almost everywhere

\textbf{Relationship to the headword:} This adverb describes how widely something exists or occurs.

\textbf{Usage note:} It is much less common than ubiquitous and is mainly used in formal writing.

\subsubsection{Examples}
\begin{enumerate}
  \item The software is now used ubiquitously across the industry.
  \item The symbol appears ubiquitously in the artist's later work.
\end{enumerate}

\section{Common collocations}
\sectionrule
\subsection{become ubiquitous}
\textbf{Applies to:} Present almost everywhere

\textbf{Meaning:} become extremely common or widely present

\textbf{Pattern:} become ubiquitous in + place or field

\textbf{Example:} Digital payment systems have become ubiquitous in many cities.

\subsection{increasingly ubiquitous}
\textbf{Applies to:} Present almost everywhere

\textbf{Meaning:} becoming present in more and more places

\textbf{Pattern:} increasingly ubiquitous in + area

\textbf{Example:} Artificial intelligence is increasingly ubiquitous in everyday products.

\subsection{nearly ubiquitous}
\textbf{Applies to:} Present almost everywhere

\textbf{Meaning:} present in almost every relevant place or case

\textbf{Example:} Internet access is nearly ubiquitous among university students.

\subsection{ubiquitous presence}
\textbf{Applies to:} Present almost everywhere

\textbf{Meaning:} the noticeable existence of something almost everywhere

\textbf{Pattern:} the ubiquitous presence of + noun

\textbf{Example:} The ubiquitous presence of advertising influences consumer behavior.

\subsection{ubiquitous technology}
\textbf{Applies to:} Present almost everywhere

\textbf{Meaning:} technology that is widely available and commonly encountered

\textbf{Example:} Ubiquitous technology has transformed the modern workplace.

\subsection{ubiquitous feature}
\textbf{Applies to:} Present almost everywhere

\textbf{Meaning:} a feature that appears in almost every relevant example or setting

\textbf{Example:} Online reviews are now a ubiquitous feature of electronic commerce.

\subsection{ubiquitous in modern life}
\textbf{Applies to:} Present almost everywhere

\textbf{Meaning:} extremely common in contemporary society

\textbf{Example:} Digital screens have become ubiquitous in modern life.

\subsection{the ubiquity of}
\textbf{Applies to:} Present almost everywhere

\textbf{Meaning:} the state of being present almost everywhere

\textbf{Pattern:} the ubiquity of + noun

\textbf{Example:} The ubiquity of mobile devices has created new privacy concerns.

\textbf{Usage note:} Ubiquity is the noun form.

\section{Synonyms and antonyms}
\sectionrule
\subsection{Close synonyms}
\subsubsection{widespread}
\textbf{Part of speech:} adjective

\textbf{Applies to:} Present almost everywhere

\textbf{Shared meaning:} Both words describe something that exists across a large area or among many people.

\textbf{Difference:} Widespread means existing in many places or affecting many people. Ubiquitous is stronger and suggests that something seems to be almost everywhere.

\textbf{Example:} The new policy received widespread public support.

\subsubsection{pervasive}
\textbf{Part of speech:} adjective

\textbf{Applies to:} Present almost everywhere

\textbf{Shared meaning:} Both words can describe something present throughout an environment.

\textbf{Difference:} Pervasive often suggests that something spreads through and influences every part of a place, system, or situation. It is frequently used for problems, attitudes, smells, or influences. Ubiquitous focuses more simply on being found almost everywhere.

\textbf{Example:} The report identified a pervasive culture of secrecy within the organization.

\subsubsection{omnipresent}
\textbf{Part of speech:} adjective

\textbf{Applies to:} Present almost everywhere

\textbf{Shared meaning:} Both words express the idea of being present everywhere or almost everywhere.

\textbf{Difference:} Omnipresent literally suggests being present everywhere and is often used for a god, authority, influence, or something that feels impossible to escape. Ubiquitous is more common in academic descriptions of widely present objects or phenomena.

\textbf{Example:} The ruler's image was omnipresent throughout the capital.

\subsubsection{common}
\textbf{Part of speech:} adjective

\textbf{Applies to:} Present almost everywhere

\textbf{Shared meaning:} Both words can describe something frequently encountered.

\textbf{Difference:} Common means frequently found or occurring. Ubiquitous is much stronger and emphasizes extremely wide presence.

\textbf{Example:} This bird is common in urban parks.

\section{Etymology}
\sectionrule
\textbf{Origin:} The word comes from Latin ubique, meaning everywhere.

\section{Practice exercises}
\sectionrule
\subsection{Word family choice}
Choose the best member of the word family for each blank.

\begin{enumerate}
  \item The \_\_\_\_\_ of digital devices has changed how students obtain information.
  \item Wireless networks are now \_\_\_\_\_ in many public buildings.
  \item The technology is used \_\_\_\_\_ across the healthcare industry.
\end{enumerate}

\subsection{Correct or incorrect?}
Decide whether each sentence is correct. Correct the inaccurate ones.

\begin{enumerate}
  \item Smartphones have become ubiquitous among young adults.
  \item The internet is ubiquitous everywhere in modern society.
  \item It is an ubiquitous problem in large cities.
  \item Security cameras are increasingly ubiquitous in public spaces.
\end{enumerate}

\subsection{Collocation practice}
Complete each sentence with a natural collocation.

\begin{enumerate}
  \item Online shopping has \_\_\_\_\_ ubiquitous in many countries.
  \item Researchers examined the ubiquitous \_\_\_\_\_ of microplastics in the environment.
  \item The \_\_\_\_\_ of social media has affected political communication.
\end{enumerate}

\subsection{Complete answer key}
\subsubsection{Word family choice}
\begin{enumerate}
  \item \textbf{ubiquity.} The ubiquity of digital devices has changed how students obtain information. \emph{Use the noun ubiquity after the and before of.}
  \item \textbf{ubiquitous.} Wireless networks are now ubiquitous in many public buildings. \emph{Use the adjective ubiquitous after are.}
  \item \textbf{ubiquitously.} The technology is used ubiquitously across the healthcare industry. \emph{Use the adverb ubiquitously to describe how the technology is used.}
\end{enumerate}

\subsubsection{Correct or incorrect?}
\begin{enumerate}
  \item \textbf{Correct} \emph{Become ubiquitous is a natural pattern for something that has become extremely common.}
  \item \textbf{Incorrect}. The internet is ubiquitous in modern society. \emph{Everywhere is unnecessary because ubiquitous already contains that meaning.}
  \item \textbf{Incorrect}. It is a ubiquitous problem in large cities. \emph{Use a because ubiquitous begins with the consonant sound /j/.}
  \item \textbf{Correct} \emph{Increasingly ubiquitous correctly describes something appearing in more and more places.}
\end{enumerate}

\subsubsection{Collocation practice}
\begin{enumerate}
  \item \textbf{become.} Online shopping has become ubiquitous in many countries. \emph{Become ubiquitous is a common collocation.}
  \item \textbf{presence.} Researchers examined the ubiquitous presence of microplastics in the environment. \emph{Ubiquitous presence describes something that can be found almost everywhere.}
  \item \textbf{ubiquity.} The ubiquity of social media has affected political communication. \emph{The ubiquity of is a common formal expression.}
\end{enumerate}

\vfill
\begin{center}
\small\color{Muted} Created and compiled by \textbf{Amir Kooshky}\\
\href{https://t.me/KooshkyTOEFL}{Telegram Channel}\quad\textbullet\quad
\href{https://www.instagram.com/kooshkytoefl}{Instagram}\quad\textbullet\quad
\href{https://t.me/KooshkyTOEFL\_pv}{Personal Telegram}
\end{center}
\end{document}
